\section{Future Work}
\label{sec:disc-future-work}

Each direction below answers a limitation named above rather than opening a new front, so the future work follows on from what the thesis has established.

The most immediate is to fix the context-preserving restore on small reflows.
A minor change, such as a single font-size step, moves the reader only slightly, but the semantic-restart strategy re-establishes the sentence near the top of the reading area and so moves them too far; a restore that simply keeps the captured offset would suit these small changes better.
More generally, replacing the hand-tuned line-bias constant with a Kalman filter would give the gaze mapping a proper estimate of its own uncertainty in place of a fixed discount, handling the same vertical jitter on a sounder basis \cite{kalman1960filtering}.

A second direction is to run the platform in automated mode at scale.
A single advisory-mode validation run has shown that the in-process decision latency and the out-of-process provider round trip can be measured end to end and stay within budget (\cref{sec:eval-performance}), but it used a rule-based reference provider on one host and gated every proposal through the researcher.
A campaign with a real decision model, an autonomous execution mode, and many sessions would turn that one-run check into a measured result and put the automated decision path under sustained measurement over a realistic deployment.

A third direction is the controlled efficacy study that the thesis deliberately left out.
Now that the platform produces reproducible, replayable records, a properly powered human-subjects study could test whether specific interventions improve reading, rather than only measuring their immediate disruption; this would turn the platform from an instrument that can run such a study into the source of its result.

A fourth direction strengthens the modularity claim from the outside.
Inviting decision and analysis providers beyond the single collaborator model would move the integration contract from author-validated toward independently validated, and adding webcam or other multimodal sources behind the existing sensing adapter would test the sensing boundary against signals quite unlike the eye tracker.
Both extend exactly the seams the architecture was built to expose (\cref{subsec:gap-sensing}, \cref{subsec:gap-integration}).
